\paragrafos{Κατακορυφήν, συμπληρωματικές και παραπληρωματικές γωνίες}{Β.3.6, «Κατακορυφήν, συμπληρωματικές και παραπληρωματικές γωνίες»}

\begin{anagke}
Δύο μονοπάτια στον κήπο της θείας Μαρίας διασταυρώνονται σαν ένα μεγάλο «Χ». Η Αγγελική μετράει με το μοιρογνωμόνιο τη μία γωνία και βρίσκει $48^\circ$. «Πόσες είναι οι άλλες τρεις;», ρωτά. Ο Νίκος μετράει όλες: $48^\circ$, $132^\circ$, $48^\circ$, $132^\circ$. «Οι απέναντι είναι ίσες! Και οι διπλανές κάνουν μαζί $180^\circ$». Η Πολυξένη ρωτά: «Είναι τυχαίο; Αν στρέψω τη μία ευθεία, θα αλλάξει;». Δοκιμάζουν με άλλη γωνία, $70^\circ$, και βρίσκουν $70^\circ$, $110^\circ$, $70^\circ$, $110^\circ$.

Η Ρήνη θέλει να φυτέψει τριαντάφυλλα σε μια γωνιά του φράχτη που είναι ορθή. Ένα παρτέρι με λουλούδια παίρνει γωνία $35^\circ$ και ρωτά πόση γωνία μένει για τα τριαντάφυλλα. Ο Στέλιος υπολογίζει: «Εδώ το άθροισμα είναι $90^\circ$, όχι $180^\circ$».

Έτσι η παρέα συναντά τρεις έννοιες: τις γωνίες που αθροίζουν $90^\circ$ (συμπληρωματικές), εκείνες που αθροίζουν $180^\circ$ (παραπληρωματικές) και αυτές που βρίσκονται απέναντι όταν δύο ευθείες τέμνονται (κατακορυφήν). Σε αυτή την παράγραφο θα δούμε και γιατί οι κατακορυφήν είναι ίσες, χωρίς μοιρογνωμόνιο.
\end{anagke}

\begin{lexi}{εφεξής, συμπληρωματική, παραπληρωματική, κατακορυφήν}{γωνίες που σχετίζονται}
Δύο γωνίες είναι \lexiword{εφεξής} όταν η μία έρχεται «από κοντά» στην άλλη. Η \lexiword{συμπληρωματική} γωνία «συμπληρώνει» μια γωνία μέχρι την ορθή. Η \lexiword{παραπληρωματική} «παρα-πληρώνει», δηλαδή συμπληρώνει ως τη μισή στροφή, τις $180^\circ$. Οι \lexiword{κατακορυφήν} γωνίες είναι αυτές που βρίσκονται «κατά κορυφήν», δηλαδή απέναντι, με κοινή κορυφή.
\end{lexi}

\begin{metafrasi}
Δύο γωνίες με κοινή κορυφή και κοινή πλευρά, χωρίς άλλο κοινό σημείο. & εφεξής\\
Δύο γωνίες με άθροισμα $90^\circ$. & συμπληρωματικές: $\hat{\alpha}+\hat{\beta}=90^\circ$\\
Δύο γωνίες με άθροισμα $180^\circ$. & παραπληρωματικές: $\hat{\alpha}+\hat{\beta}=180^\circ$\\
Οι απέναντι γωνίες δύο τεμνόμενων ευθειών. & κατακορυφήν: $\hat{\alpha}=\hat{\gamma}$\\
Η συμπληρωματική της $35^\circ$. & $90^\circ-35^\circ=55^\circ$\\
Η παραπληρωματική της $130^\circ$. & $180^\circ-130^\circ=50^\circ$\\
\end{metafrasi}

\tmima{Η θεωρία}

\tmimasmall{Εφεξής και διαδοχικές γωνίες}

\begin{orismos}{εφεξής γωνίες}
Δύο γωνίες λέγονται \textbf{εφεξής} όταν έχουν κοινή κορυφή, μία κοινή πλευρά και κανένα άλλο κοινό σημείο.
\end{orismos}

Μια σειρά από γωνίες που η καθεμία είναι εφεξής με την επόμενη λέγονται \textbf{διαδοχικές}. Για παράδειγμα στο σχήμα οι γωνίες $\mathrm{ΑΟΒ}$, $\mathrm{ΒΟΓ}$ και $\mathrm{ΓΟΔ}$ είναι διαδοχικές. Το άθροισμά τους είναι η γωνία $\mathrm{ΑΟΔ}$.

\begin{center}
\begin{tikzpicture}[scale=0.9]
\draw[thick] (0,0)--(4.5,0) node[right] {$\mathrm{Α}$};
\draw[thick] (0,0)--({4.2*cos(35)},{4.2*sin(35)}) node[above right] {$\mathrm{Β}$};
\draw[thick] (0,0)--({4.2*cos(80)},{4.2*sin(80)}) node[above] {$\Gamma$};
\draw[thick] (0,0)--({4.2*cos(130)},{4.2*sin(130)}) node[above left] {$\Delta$};
\fill (0,0) circle (2pt) node[below left] {$\mathrm{Ο}$};
\draw[gblue] (1.0,0) arc(0:35:1.0); \draw[gblue] (1.2,0.55) arc(25:80:1.2); \draw[gblue] (0.7,0.5) arc(35:130:0.7);
\end{tikzpicture}
\end{center}

\tmimasmall{Συμπληρωματικές και παραπληρωματικές}

\begin{orismos}{συμπληρωματικές και παραπληρωματικές}
Δύο γωνίες λέγονται \textbf{συμπληρωματικές} όταν το άθροισμά τους είναι $90^\circ$ και \textbf{παραπληρωματικές} όταν το άθροισμά τους είναι $180^\circ$. Κάθε μία από τις δύο λέγεται συμπληρωματική (ή παραπληρωματική) της άλλης.
\end{orismos}

Για παράδειγμα η συμπληρωματική της $35^\circ$ είναι η $55^\circ$ και η παραπληρωματική της $130^\circ$ είναι η $50^\circ$. Δύο γωνίες δεν χρειάζεται να είναι εφεξής για να είναι συμπληρωματικές ή παραπληρωματικές: αρκεί το άθροισμα των μέτρων τους.

\begin{center}
\begin{tikzpicture}[scale=0.8]
\draw[thick] (-0.3,0)--(3.5,0); \draw[thick] (0,-0.3)--(0,3.5); \draw[thick] (0,0)--({3.3*cos(35)},{3.3*sin(35)});
\draw[gblue] (0.8,0) arc(0:35:0.8); \draw[gwarmD] (0.6,0.4) arc(35:90:0.6); \draw (0.3,0)--(0.3,0.3)--(0,0.3);
\node[gblue] at (1.2,0.3) {\small $\hat{\alpha}$}; \node[gwarmD] at (0.4,1.1) {\small $\hat{\beta}$}; \node at (1.7,-0.7) {\small $\hat{\alpha}+\hat{\beta}=90^\circ$};
\begin{scope}[xshift=6.2cm]
\draw[thick] (-2.2,0)--(2.6,0); \draw[thick] (0,0)--({2.3*cos(60)},{2.3*sin(60)});
\draw[gblue] (0.7,0) arc(0:60:0.7); \draw[gwarmD] (-0.55,0) arc(180:60:0.55);
\node[gblue] at (1.05,0.3) {\small $\hat{\alpha}$}; \node[gwarmD] at (-0.3,0.75) {\small $\hat{\beta}$}; \node at (0.3,-0.7) {\small $\hat{\alpha}+\hat{\beta}=180^\circ$};
\end{scope}
\end{tikzpicture}
\end{center}

\tmimasmall{Κατακορυφήν γωνίες}

\begin{orismos}{κατακορυφήν γωνίες}
Δύο γωνίες λέγονται \textbf{κατακορυφήν} όταν έχουν κοινή κορυφή και οι πλευρές τους είναι αντικείμενες ημιευθείες.
\end{orismos}

Όταν δύο ευθείες τέμνονται σχηματίζονται τέσσερις γωνίες, και οι δύο ζεύγη απέναντι γωνιών είναι κατακορυφήν.

\begin{center}
\begin{tikzpicture}[scale=1.1]
\draw[thick,{Stealth}-{Stealth}] (-2.5,-1)--(2.5,1); \draw[thick,{Stealth}-{Stealth}] (-1.2,-1.8)--(1.2,1.8);
\fill (0,0) circle (2pt); \node[below right] at (0.05,0) {$\mathrm{Ο}$};
\node[gblue] at (1.2,0.15) {$\hat{\alpha}$}; \node[gblue] at (-1.2,-0.15) {$\hat{\gamma}$}; \node at (0.2,0.9) {$\hat{\beta}$}; \node at (-0.2,-0.9) {$\hat{\delta}$};
\end{tikzpicture}
\end{center}

\begin{idiotita}{κατακορυφήν γωνίες}
Οι κατακορυφήν γωνίες είναι \textbf{ίσες}.
\end{idiotita}

Γιατί: Η $\hat{\alpha}$ και η $\hat{\beta}$ είναι παραπληρωματικές ($\hat{\alpha}+\hat{\beta}=180^\circ$), και η $\hat{\gamma}$ και η $\hat{\beta}$ επίσης ($\hat{\gamma}+\hat{\beta}=180^\circ$). Άρα $\hat{\alpha}=180^\circ-\hat{\beta}=\hat{\gamma}$.

\begin{prosoxi}
\begin{itemize}[leftmargin=1.2em,itemsep=0.3ex,topsep=2pt]
\item Συμπληρωματικές: $90^\circ$. Παραπληρωματικές: $180^\circ$. Η «παραπληρωματική» είναι η μεγαλύτερη λέξη και το μεγαλύτερο άθροισμα.
\item Δεν υπάρχει συμπληρωματική γωνία για γωνία ίση ή μεγαλύτερη από $90^\circ$ (θα βγαίνε μηδέν ή αρνητικό).
\item Οι διπλανές γωνίες δύο τεμνόμενων ευθειών δεν είναι ίσες (εκτός αν είναι ορθές): είναι παραπληρωματικές.
\item Δύο ίσες γωνίες δεν είναι πάντα κατακορυφήν: αρκεί να έχουν ίδιο μέτρο.
\end{itemize}
\end{prosoxi}

\tmima{Λυμένα παραδείγματα}

\paradeigma{Συμπλήρωμα και παραπλήρωμα}
\tixero Οι γωνίες $35^\circ$ και $130^\circ$.\par
\tithelo Η συμπληρωματική της πρώτης και η παραπληρωματική της δεύτερης.\par
\possindeo Αφαιρώ από το $90^\circ$ και από το $180^\circ$ αντίστοιχα.\par
\apantisi $90^\circ-35^\circ=55^\circ$ και $180^\circ-130^\circ=50^\circ$.\par
\elegxos $35^\circ+55^\circ=90^\circ$ και $130^\circ+50^\circ=180^\circ$.

\paradeigma{Δύο τεμνόμενες ευθείες}
\tixero Δύο ευθείες τέμνονται και μία από τις γωνίες είναι $50^\circ$.\par
\tithelo Οι άλλες τρεις γωνίες.\par
\possindeo Η απέναντί της είναι κατακορυφήν, άρα $50^\circ$. Οι διπλανές της είναι παραπληρωματικές της, δηλαδή $180^\circ-50^\circ=130^\circ$, και η τέταρτη είναι κατακορυφήν της μίας από αυτές, άρα επίσης $130^\circ$.\par
\apantisi $50^\circ$, $130^\circ$ και $130^\circ$.\par
\elegxos Το άθροισμα και των τεσσάρων: $50^\circ+130^\circ+50^\circ+130^\circ=360^\circ$.

\paradeigma{Η συμπληρωματική με το διπλάσιο}
\tixero Μια γωνία είναι συμπληρωματική με μια άλλη που είναι διπλάσια από αυτήν.\par
\tithelo Οι δύο γωνίες.\par
\possindeo Αν η μία είναι $\omega$, η άλλη είναι $2\omega$ και το άθροισμά τους είναι $90^\circ$: $3\omega=90^\circ$, άρα $\omega=30^\circ$.\par
\apantisi $30^\circ$ και $60^\circ$.\par
\elegxos $30^\circ+60^\circ=90^\circ$ και η δεύτερη είναι διπλάσια της πρώτης.

\paradeigma{Όλες οι κατακορυφήν}
\tixero Δύο ευθείες τέμνονται και μία γωνία είναι $\hat{\alpha}=2x+10^\circ$ ... \par
\tithelo ...\par
\possindeo ...\par
\apantisi ...\par
\elegxos ...

\tmima{Ασκήσεις}

\zoni{Το κατάλαβα}{συμπληρωματικές, παραπληρωματικές, κατακορυφήν}

\ex Σωστό ή λάθος; α) Δύο γωνίες με άθροισμα $90^\circ$ λέγονται παραπληρωματικές. \quad β) Δύο κατακορυφήν γωνίες είναι πάντα ίσες. \quad γ) Η συμπληρωματική της μηδενικής γωνίας είναι η ορθή. \quad δ) Η παραπληρωματική μιας αμβλείας είναι πάντα οξεία. \quad ε) Δύο ίσες γωνίες είναι πάντα κατακορυφήν. \quad στ) Δύο εφεξής γωνίες έχουν κοινή κορυφή και μία κοινή πλευρά.

\ex Βρες τη συμπληρωματική των γωνιών: $35^\circ$, $62^\circ$, $8^\circ$, $89^\circ$.

\ex Βρες την παραπληρωματική των γωνιών: $130^\circ$, $90^\circ$, $5^\circ$, $111^\circ$.

\ex Δύο ευθείες τέμνονται και μία γωνία είναι $48^\circ$. Βρες τις άλλες τρεις γωνίες.

\ex Δύο ευθείες $x\mathrm{Ο}x'$ και $y\mathrm{Ο}y'$ τέμνονται στο $\mathrm{Ο}$. Γράψε όλα τα ζεύγη κατακορυφήν γωνιών και όλα τα ζεύγη παραπληρωματικών γωνιών που σχηματίζονται.

\ex Σχεδίασε μια γωνία $40^\circ$ και μετά τη συμπληρωματική της και την παραπληρωματική της.

\ex Τι είδους γωνία είναι: α) η παραπληρωματική μιας οξείας, β) η παραπληρωματική μιας ορθής, γ) η παραπληρωματική μιας αμβλείας, δ) η συμπληρωματική μιας οξείας;

\ex Οι γωνίες $\mathrm{ΑΟΒ}=30^\circ$, $\mathrm{ΒΟΓ}=45^\circ$ και $\mathrm{ΓΟΔ}=50^\circ$ είναι διαδοχικές. Πόση είναι η γωνία $\mathrm{ΑΟΔ}$; Είναι οξεία, ορθή ή αμβλεία;

\zoni{Το χρησιμοποιώ}{διασταυρώσεις, σκάλες, ψαλίδια}

\ex \textbf{Η διασταύρωση.} Δύο μονοπάτια του κήπου τέμνονται και μία γωνία είναι $70^\circ$. Πόσες μοίρες έχουν οι άλλες γωνίες; Πόσες διαφορετικές τιμές εμφανίζονται;

\ex \textbf{Τα τριαντάφυλλα.} Στη γωνιά του φράχτη, που είναι ορθή, φυτεύεται ένα παρτέρι με γωνία $35^\circ$ και το υπόλοιπο φυτεύεται με τριαντάφυλλα. Πόση γωνία μένει για τα τριαντάφυλλα;

\ex \textbf{Η σκάλα.} Μια σκάλα ακουμπά σε έναν τοίχο και σχηματίζει με το έδαφος γωνία $72^\circ$. Πόση γωνία σχηματίζει με τον τοίχο; (Ο τοίχος είναι κάθετος στο έδαφος.)

\ex \textbf{Το ψαλίδι.} Ο παππούς Γιάννης ανοίγει ένα ψαλίδι ώστε οι δύο λεπίδες να σχηματίζουν γωνία $40^\circ$. Πόση γωνία σχηματίζουν οι λαβές του;

\ex \textbf{Η πίτσα.} Μια γωνία είναι ίση με τα $\frac{2}{5}$ της ορθής. Πόσες μοίρες είναι; Βρες τη συμπληρωματική και την παραπληρωματική της.

\zoni{Το σκέφτομαι}{γιατί και αν}

\ex Εξήγησε γιατί οι κατακορυφήν γωνίες είναι ίσες, χρησιμοποιώντας τις παραπληρωματικές γωνίες.

\ex Δύο ίσες γωνίες είναι απαραίτητα κατακορυφήν; Δώσε ένα αντιπαράδειγμα.

\ex Βρες το μέτρο της γωνίας για την οποία ισχύει ότι: α) η συμπληρωματική της είναι κατά $20^\circ$ μεγαλύτερη από αυτήν, β) η παραπληρωματική της είναι κατά $40^\circ$ μικρότερη από αυτήν, γ) η συμπληρωματική της είναι διπλάσια από αυτήν, δ) η παραπληρωματική της είναι τριπλάσια από αυτήν.

\ex Υπάρχει γωνία ίση με τη συμπληρωματική της; Και γωνία ίση με την παραπληρωματική της; Βρες τες.

\zoni{Για το σπίτι}{ψάξε Χ}

\ex Σχεδίασε δύο τεμνόμενες ευθείες, μέτρησε και τις τέσσερις γωνίες και έλεγξε ότι οι κατακορυφήν είναι ίσες.

\ex Βρες στο σπίτι ένα ζευγάρι τεμνόμενων ευθειών (ψαλίδι, παράθυρο, αναδιπλούμενη καρέκλα) και μέτρησε τις γωνίες που σχηματίζουν.
